%preamble: defining document parameters and packages

\documentclass[11pt]{report}         %defining documentclass
\usepackage[brazilian] {babel}          %defining document language
\usepackage[utf8]{inputenc}          %defining to support portuguese a accentuation
\usepackage{authblk}                 %defining packages to better author format
\usepackage{indentfirst}             %defining packages to indention first paragraph's line
\usepackage[margin=2.5cm]{geometry}  %defining margin
\usepackage[table]{xcolor}           %defining table colors
\usepackage{pdfpages}                %enabling pdf attachment
\usepackage{graphicx}                %enabling image attachment
\usepackage{float}                   %enabling a better abaut image em tabular posicions in pages

%cover (title, author and date)

\title{\textbf{Relatório da AP3}}
\author{Arthur Siqueira Paz Teixeira \\ \& \\ Yure da Silva Lacerda de Barros} 
\affil{Disciplina: Circuitos Lógicos -- Professor: Fabio Marujo}
\date{\today}


\begin{document}                  %start document

\maketitle                        %generate the cover
\tableofcontents                  %generate the summary

\chapter {Introdução}             %First chapter

Neste documento está descrito o planejamento, conforme pedido no anexo \ref{anexo1}, e execução do circuito somador BCD de dois números (de digito único) 
com saida simples em um unico dysplay de LED de 7 segmentos. Para que seja possivel a extecução da representação de resultados acima de 9 
(limitados entre 10 e 18) será usado, recursivamente, um LED para representar a presensa ou ausencia do número 1 na casa das dezenas.

Além disso também foi proposto que projete-se o módulo básico de um comparador de dois números de 4 bits, expansível utilizando a ferramenta MACRO do 
Circuit Maker para implementar o módulo básico na simulação para bits.


\chapter {Planejamento}

Como parte do planejamento deve-se pontuar a divisão do processo em 5 partes, definir entradas e saidas, montagem da tabela verdade, simplificação do
circuito, simulação e lista de componetes necessarios (no caso do meio somador que foi implementado de forma prática).

\section {Definindo entradas e saidas}

\begin {itemize}

\item \textbf{Somador BCD:} Precisa-se de dois C.I.s somadores para o projeto, cada somador contem um total de oito entradas: $A1_{in}$, $A2_{in}$, $A3_{in}$, $A4_{in}$, $B1_{in}$, 
$B2_{in}$, $B3_{in}$ e $B4_{in}$. Essas entradas geram quatro saídas: $S3_{out}$, $S2_{out}$, $S1_{out}$, $S0_{out}$ e $C_{out}$. Para o resultado pedido 
neste projeto as saidas do primeiro C.I. serão conectadas nas entradas $A1_{in}$, $A2_{in}$, $A3_{in}$ e $A4_{in}$ do segundo somador, sobrando apenas 
4 entradas que serão conectadas de forma a auxiliar em problemas que vermos a seguir.

Na soma de BCD, só são aceitos resustados que vão de 0 (0000) a 9 (1001). Quando você soma dois dígitos BCD usando um somador binário comum (o primeiro 7483), 
duas questões aparecem se o resultado for maior que 9:

\begin {itemize}

\item \textbf{Resultados Inválidos:} A soma binária vai resultar valores de 10 (1010) a 15 (1111), que não existem no código BCD e não podem 
ser representados em um único o display de 7 segmentos.

\item \textbf{Vai-Um(textit{Carry in}):} Quando a soma der entre 16 (10000) e 19 (10011), o somador, que possui apenas 4 bits de saída, precisa de um bit extra 
e joga esse "vai-um" na saida $C_{out}$ para 1. Para o BCD, 16 deveria ser o display de dezenas em 1 e o de unidades em 6 (0110), e não 0 (0000).

\end {itemize}

A solução é considerar uma defasagem de 6 digitos sempre que o valor for de 9 a 19. Considerando isso o segundo display deverá somar 6 em todos os resultados 
dentro destas caracteristicas.

\item \textbf{Comparador:} O circuito faz a analise da entrada a cada bit, do mais significativo(MSB) até o menos(LSB), para verificar se A é igual, maior 
ou menor que B.

\begin{itemize}

\item \textbf{Definiu-se:} 

\begin {itemize} 

\item $A = [A_3, A_2, A_1, A_0]$ e $B = [B_3, B_2, B_1, B_0]$ onde $X_3$ é o MSB.

\item para igualdade todos os bits devem ser iguais: $[A_3, A_2, A_1, A_0] = [B_3, B_2, B_1, B_0]$

\item para maioridade: Definiu-se uma hierarquia onde se $A_3 > B_3$. Logo, A já é maior que B, pois o MSB de A é maior que o de B. Está relação se repete 
em cada um dos bits comparados

\item para minoridade: Definiu-se o inverso do item anterior, se $A_3 < B_3$. Logo, A já é menor que B, pois o MSB de A é menor que o de B.

\item para expandir: Coloca-se uma porta OR em cada uma das saídas. A primeira entrada da porta OR é a saída correspondente (X ou Y) e a segunda 
é a saída correspondente (X ou Y) do comparador anterior. Assim, consegue-se expandir em 8 bits a cada comparador anexado.

\end {itemize} 
\end {itemize}
\end {itemize}

\section {Montagem da tabela verdade}

Enquanto esta etapa é dispensvel no comparador, que tem por objetivo ser expansível e, portanto não deve ser simplificado para evitar problemas 
na expanção, a tabela do somador deve focar em nos resultados entre 9 e 19 para que possamos montar o circuito de forma que a soma de 6 somente ocorra 
quando necessaria.

\begin{table}[h!]
    \centering
    
    \caption{Tabela Verdade Teórica/Simulada}
    
    \begin{tabular}{|c|c|c|c|c|c|>{\columncolor{black}\color{white}}c|}

      \hline
      \textbf{Resultado da Soma} & \textbf{C} & \textbf{S3} & \textbf{S2} & \textbf{S1} & \textbf{S0} & \textbf{X} \\ \hline
      0  & 0 & 0 & 0 & 0 & 0 & 0                  \\ \hline
      1 até 7 & ... & ... & ... & ... & ... & 0   \\ \hline
      8  & 0 & 1 & 0 & 0 & 0 & 0                  \\ \hline
      9  & 0 & 1 & 0 & 0 & 1 & 0                  \\ \hline
      10 & 0 & 1 & 0 & 1 & 0 & 1                  \\ \hline
      11 & 0 & 1 & 0 & 1 & 1 & 1                  \\ \hline
      12 & 0 & 1 & 1 & 0 & 0 & 1                  \\ \hline
      13 & 0 & 1 & 1 & 0 & 1 & 1                  \\ \hline
      14 & 0 & 1 & 1 & 1 & 0 & 1                  \\ \hline
      15 & 0 & 1 & 1 & 1 & 1 & 1                  \\ \hline
      16 & 1 & 0 & 0 & 0 & 0 & 1                  \\ \hline
      17 & 1 & 0 & 0 & 0 & 1 & 1                  \\ \hline
      18 & 1 & 0 & 0 & 1 & 0 & 1                  \\ \hline
      19 & 1 & 0 & 0 & 1 & 1 & 1                  \\ \hline
      20 até 31 & ... & ... & ... & ... & ... & X \\ \hline

    \end{tabular}

\end{table}

.\\

\section {Simplificação do circuito}

Extraindo da tabela verdade os valores que são relevantes para a simplificação temos:


\begin{itemize}

\item \textbf{mapeando as posições na tabela}


\begin{table}[h!]
    \centering
    
    \caption{Tabela mapeada}
    
    \begin{tabular}{|c|c|c|c|c|}

      \hline
      \textbf{---} & \textbf{S1' S0'} & \textbf{S1' S0} & \textbf{S1 S0} & \textbf{S1 S0'}\\ \hline
      C' S3' S2' & 0 & 1 & 3 & 2       \\ \hline
      C' S3' S2  & 4 & 5 & 7 & 6       \\ \hline
      C' S3 S2   & 12 & 13 & 15 & 14   \\ \hline
      C' S3 S2'  & 8 & 9 & 11 & 10     \\ \hline
      C S3' S2'  & 16 & 17 & 19 & 18   \\ \hline
      C S3' S2   & 20 & 21 & 23 & 22   \\ \hline
      C S3 S2    & 28 & 29 & 31 & 30   \\ \hline
      C S3 S2'   & 24 & 25 & 27 & 26   \\ \hline
      
    \end{tabular}

\end{table}

\newpage

\item \textbf{Mapa de Karnaugh preenchido}

\begin{table}[h!]
    \centering
    
    \caption{Mapa de Karnaugh}
    
    \begin{tabular}{|c|c|c|c|c|}

      \hline
      \textbf{---} & \textbf{S1' S0'} & \textbf{S1' S0} & \textbf{S1 S0} & \textbf{S1 S0'}\\ \hline
      C' S3' S2' & 0 & 0 & 0 & 0 \\ \hline
      C' S3' S2  & 0 & 0 & 0 & 0 \\ \hline
      C' S3 S2   & 1 & 1 & 1 & 1 \\ \hline
      C' S3 S2'  & 0 & 0 & 1 & 1 \\ \hline
      C S3' S2'  & 1 & 1 & 1 & 1 \\ \hline
      C S3' S2   & X & X & X & X \\ \hline
      C S3 S2    & X & X & X & X \\ \hline
      C S3 S2'   & X & X & X & X \\ \hline
      
    \end{tabular}

\end{table}

.\\

\item \textbf{Equação encontrada}

$X = C + (S3 \cdot S1) + (S3 \cdot S2)$

\end{itemize}

\section{Simulação}

\begin{itemize}

\item \textbf{Somador:} Para verificar o funcionamento antes da montagem em um circuito fisico é necessaria a implementação do circuito em um softwere que simule o seu 
funcionamento, evitando gastos desnecessarios e perca de tempo na correção de possiveis erros. A imagem a seguir demonstra a implementação do Circuito 
desenvolvido neste documento

\begin{figure}[H] % <--- [H] força a imagem a ficar exatamente AQUI no texto
    \centering
    \includegraphics[width=0.7\textwidth]{/home/sparthur/git/logic-circuits-design/AP3/circuit-schematic.png}
    \caption{Diagrama Esquemático do Circuito}
    \label{fig:esquematico}
\end{figure}

nesta simulação também foram feitas analises das entradas e saidas, montando uma tabela verdade

\newpage

\textbf{Tabela Simulada}

\begin{table}[h!]
    \begingroup
    \footnotesize
    \caption{Tabela Simulada}

        \begin{tabular}{|>{\columncolor{black}\color{white}}c|c|c|c|c|c|c|c|c|c|}

      \hline
      \textbf{Resultado} & \textbf{Entrada A} & \textbf{Entrada B} & \textbf{Saída S} & \textbf{C} & \textbf{X} & \textbf{Entrada S2} & \textbf{Saída S2} & \textbf{LED} & \textbf{7 segmentos}\\ \hline
      00 & 0000 (0) & 0000 (0) & 0000 (00) & 0 & 0 & 0000 (0) & 0000 (0) & 0 & 0 \\ \hline
      01 & 0001 (1) & 0000 (0) & 0001 (01) & 0 & 0 & 0000 (0) & 0001 (1) & 0 & 1 \\ \hline
      02 & 0001 (1) & 0001 (1) & 0010 (02) & 0 & 0 & 0000 (0) & 0010 (2) & 0 & 2 \\ \hline
      03 & 0010 (2) & 0001 (1) & 0011 (03) & 0 & 0 & 0000 (0) & 0011 (3) & 0 & 3 \\ \hline
      04 & 0010 (2) & 0010 (2) & 0100 (04) & 0 & 0 & 0000 (0) & 0100 (4) & 0 & 4 \\ \hline
      05 & 0011 (3) & 0010 (2) & 0101 (05) & 0 & 0 & 0000 (0) & 0101 (5) & 0 & 5 \\ \hline
      06 & 0011 (3) & 0011 (3) & 0110 (06) & 0 & 0 & 0000 (0) & 0110 (6) & 0 & 6 \\ \hline
      07 & 0100 (4) & 0011 (3) & 0111 (07) & 0 & 0 & 0000 (0) & 0111 (7) & 0 & 7 \\ \hline
      08 & 0100 (4) & 0100 (4) & 1000 (08) & 0 & 0 & 0000 (0) & 1000 (8) & 0 & 8 \\ \hline
      09 & 0101 (5) & 0100 (4) & 1001 (09) & 0 & 0 & 0000 (0) & 1001 (9) & 0 & 9 \\ \hline
      10 & 0101 (5) & 0101 (5) & 1010 (10) & 0 & 1 & 0110 (6) & 0000 (0) & 1 & 0 \\ \hline
      11 & 0110 (6) & 0101 (5) & 1011 (11) & 0 & 1 & 0110 (6) & 0001 (1) & 1 & 1 \\ \hline
      12 & 0110 (6) & 0110 (6) & 1100 (12) & 0 & 1 & 0110 (6) & 0010 (2) & 1 & 2 \\ \hline
      13 & 0111 (7) & 0110 (6) & 1101 (13) & 0 & 1 & 0110 (6) & 0011 (3) & 1 & 3 \\ \hline
      14 & 0111 (7) & 0111 (7) & 1110 (14) & 0 & 1 & 0110 (6) & 0100 (4) & 1 & 4 \\ \hline
      15 & 1000 (8) & 0111 (7) & 1111 (15) & 0 & 1 & 0110 (6) & 0101 (5) & 1 & 5 \\ \hline
      16 & 1000 (8) & 1000 (8) & 0000 (00) & 1 & 1 & 0110 (6) & 0110 (6) & 1 & 6 \\ \hline
      17 & 1001 (9) & 1000 (9) & 0001 (01) & 1 & 1 & 0110 (6) & 0111 (7) & 1 & 7 \\ \hline
      18 & 1001 (9) & 1001 (9) & 0010 (02) & 1 & 1 & 0110 (6) & 1000 (8) & 1 & 8 \\ \hline

    \end{tabular}
    \endgroup
\end{table}

\newpage

\item \textbf {Comparador}

No caso do comarador foram feitas variadas simulações para cada modo de implementação:

\begin {itemize}

\item \textbf {Diagrama Esquemático do Circuito Comparador}

\begin{figure}[H] % <--- [H] força a imagem a ficar exatamente AQUI no texto
    \centering
    \includegraphics[width=0.7\textwidth]{/home/sparthur/git/logic-circuits-design/AP3/final-report/circuito-do-comparador.png}
    \caption{Diagrama Esquemático do Circuito Comparador}
    \label{fig:esquematico}
\end{figure}

\item \textbf{Diagrama Esquemático do Circuito com o comparador expansível}

\begin{figure}[H] % <--- [H] força a imagem a ficar exatamente AQUI no texto
    \centering
    \includegraphics[width=0.7\textwidth]{/home/sparthur/git/logic-circuits-design/AP3/final-report/comparador-expansível.png}
    \caption{Diagrama Esquemático do Circuito com o comparador expansível}
    \label{fig:esquematico}
\end{figure}

\newpage

\item \textbf{Diagrama Esquemático do Circuito com o comparador com $A = B$}

\begin{figure}[h] % <--- [h] força a imagem a ficar exatamente aqui no texto
    \centering
    \includegraphics[width=0.7\textwidth]{/home/sparthur/git/logic-circuits-design/AP3/final-report/comparador-no-somador-BCD-com-A-igual-A-B.png}
    \caption{Diagrama Esquemático do Circuito com o comparador com $A = B$}
    \label{fig:esquematico}
\end{figure}

\item \textbf {Diagrama Esquemático do Circuito com o comparador com $A > B$}

\begin{figure}[h] % <--- [h] força a imagem a ficar exatamente aqui no texto
    \centering
    \includegraphics[width=0.7\textwidth]{/home/sparthur/git/logic-circuits-design/AP3/final-report/comparador-no-somador-BCD-com-A-maior-que-B.png}
    \caption{Diagrama Esquemático do Circuito com o comparador com $A > B$}
    \label{fig:esquematico}
\end{figure}

\newpage

\item \textbf{Diagrama Esquemático do Circuito com o comparador com $A < B$}

\begin{figure}[h] % <--- [h] força a imagem a ficar exatamente aqui no texto
    \centering
    \includegraphics[width=0.7\textwidth]{/home/sparthur/git/logic-circuits-design/AP3/final-report/comparador-no-somador-BCD-com-B-maior-que-A.png}
    \caption{Diagrama Esquemático do Circuito com o comparador com $A < B$}
    \label{fig:esquematico}
\end{figure}

\end {itemize}
\end {itemize}

\section{Lista de componetes}

Para a aplicação serão necessarios:

\begin{itemize}

\item \textbf{Protoboard [02x]}
\item \textbf{C.I. 7483 (meio-somador) [02x]}
\item \textbf{C.I. 7448 (conversor) [01x]}
\item \textbf{C.I. 7408 (AND) [01x]}
\item \textbf{C.I. 7432 (OR) [01x]}
\item \textbf{LED [09x]}
\item \textbf{Display de 7 segmentos (Anodo Comum) [01x]}
\item \textbf{Resistores (330 Ohms) [16x]}

\end {itemize}

\chapter{Conclusão}

Ao fim da simulação ouve a montagem do circuito somador, que seguiu o diagrama projetado no \textit{Circuit-Maker-2000} e apresentado acima.
Após a montagem no \textit{Protoboard} o foram feitas comparações com os a tabela verdade e foi observada a confirmação dos dados aqui póstos 
durante a preparação.

Potanto pode-se definir como bem sucedida a prática de forma que seus principais objetivos foram compridos e um somador funcional foi montado.

\chapter {Anexos}

\includepdf[pages= 1, addtotoc={1, section, 2, {Anexo 1: Guia da AP3}, anexo1}]{/home/sparthur/git/logic-circuits-design/AP3/AP3-exercice-guide.pdf}

\end{document}
